%\documentclass[linenumbers,twocolumn]{aastex7}
\documentclass[onecolumn]{aastex7}

\newcommand{\Header}[1]{{\large{\bf{#1}}}}


\begin{document}

\title{Thoughts on research}
\author[0000-0002-6134-8946]{Ilya Mandel}
\affiliation{School of Physics and Astronomy, Monash University, Clayton VIC 3800, Australia}
\affiliation{OzGrav: The ARC Centre of Excellence for Gravitational Wave Discovery, Australia}
\email{ilya.mandel@monash.edu}


\begin{abstract}
These were written in 2007, in my last year of grad school, just after I read Cole Miller's \href{https://pages.astro.umd.edu/~mcmiller/teaching/research.pdf}{Thoughts about research} (which I highly recommend!).  At that time, having wished that I read Cole's thoughts five years previously, I wrote down the following list of purely practical suggestions based on my own failings during my PhD.  I am intentionally leaving them as is (2025), modulo light proofreading.  I would probably be more ambitious about item 1 today, but still agree with most of the rest. 
\end{abstract}

\maketitle

\vspace{0.1in}

\Header{1. Make sure your first project is well-defined and interesting}

Don't be afraid to say no to your adviser's first suggestion for a project.  Don't be afraid to say no even if you've already accepted a project, but realise that the project is going nowhere or you are not excited by it.  Your adviser should be able to help you pick a project that has a clear set of goals achievable on a reasonable time frame.  Yes, eventually you may choose some open-ended projects where the expected result is not immediately clear when you start your research, but these might not be a good choice for your first project.  It's very inefficient, not to mention discouraging, to spend several years on an unappealing task without an end in sight.

\Header{2. Working in a group vs.~working alone}

You can work on a research project on your own (with occasional advice), in a small group of two or three people, or in a large group. There are positives and negatives to all three approaches.  For example, working in a large group guarantees that you always have people very familiar with the project who have a vested interest in helping out if you are stuck, but it could also delay things when your collaborators are busy with other work.  Of course, different projects might be more or less amenable to working alone vs.~in a group, and much depends on your personal preference.  Try to work with different collaborators to see what you like best.

\Header{3. Time in the office}

Spend reasonable amounts of time in the office, and at reasonable hours. It's useful for interacting with colleagues, for making sure others remember that you exist, and for your general mental health.

\Header{4. A paper a day keeps the angry committee members away}

It's very embarrassing to realise when fielding questions after a talk or defending your candidacy that you are not familiar with the literature in your field.  It's also bad for your research.  Reading literature is an acquired skill: how do you quickly go through a paper to pull out the significant results with a critical attitude?  Perhaps a plan of reading a paper a day might help.  It could be something related to your current project, one of the historically famous papers in your field, or just an interesting paper you noticed on the archives that will help you broaden your horizons.

\Header{5. Keep good notes}

When working on problem sets, you've probably made scratch marks on the back of envelopes that you've then thrown away (perhaps along with the problem sets themselves, after they've been graded).  This is not a good approach to research notes.  You should find the most convenient way to keep them organised.  Whether you prefer to write everything in one big notebook, to keep several folders or binders for notes on different projects, or to type things up in LaTeX the moment you get an idea is up
to you.  But the important thing is to have access to your ideas and calculations from several months ago.  It might also be nice to keep notes and bibliographical information on the papers you've read, perhaps in a large .bib file so that you can easily reference them once you go to write up your results.  And if you obtain a result that you suspect will go into a future paper, you may find it helpful to write it up then and there: one reason some people have trouble getting papers out is that even if the initial research was fun, writing it up months later, after you've moved on to other projects and the results are half-forgotten, is not.

\end{document}